\chapter{Gastric Emptying Tests}
\section*{What Is This Test?} Gastric emptying tests measure how quickly food or liquid leaves the stomach.\section*{Alternative Names} Gastric emptying scan; gastric emptying scintigraphy; breath test.\section*{Principle} A labeled meal or breath marker is tracked over time to estimate stomach emptying.\section*{Why This Test Is Ordered} It evaluates nausea, vomiting, early fullness, bloating, or suspected gastroparesis and rapid emptying.\section*{Primary vs Secondary Test} It is a secondary physiologic test after history and initial evaluation.\section*{How This Test Is Performed} The patient eats a standardized meal or drinks a test substrate, with images or breath samples collected at scheduled intervals.\section*{What Preparation Is Needed} Fasting is usually required; medicines affecting motility may need adjustment only under medical direction. Report pregnancy possibility.\section*{Understanding Results} Delayed or rapid emptying must be interpreted with symptoms, meal composition, glucose, and protocol.\section*{Follow-up Tests} Endoscopy, glucose testing, imaging, or motility consultation may follow.\section*{References} \begin{enumerate}\item MedlinePlus. Gastric Emptying Tests.\item American College of Gastroenterology. Gastroparesis guidance.\end{enumerate}\clearpage\section*{Understanding Results and Follow-up}
The laboratory report should identify the specimen, method, units, reference interval or decision limit, and whether the result is positive, negative, abnormal, or inconclusive. Reference intervals vary with age, sex, pregnancy, specimen type, assay, and laboratory; a result outside a range is not automatically a diagnosis, and a result within a range does not always exclude disease. Discuss unexpected results with the ordering clinician, who can decide whether repeat or confirmatory testing, treatment, or referral is appropriate. Do not change medicines or delay urgent care based on an isolated result.
\section*{Regulatory Status and Authorization Date}This chapter describes a test category, not a specific commercial product. FDA, CDSCO, EU IVDR, and MHRA approval or registration dates therefore cannot be assigned without the assay manufacturer, product name, instrument, and intended use. For laboratory-developed tests, an individual agency approval date may not exist. See the Preface for the interpretation of regulatory status.
\clearpage
